% !TEX program = xelatex
% 本审计现归档于 src/model 对应的唯一文档目录；它是审计材料，不替代 model 模块技术手册。
% 编译方式：xelatex component_models_math_audit.tex（建议两遍以生成目录）
% src/model 数学模型工业级审计报告（中文）
\documentclass[11pt,a4paper,fontset=fandol]{ctexart}

\usepackage[margin=2.2cm]{geometry}
\usepackage{amsmath,amssymb,bm,mathtools}
\usepackage{booktabs,longtable,array,tabularx}
\usepackage[table]{xcolor}
\usepackage{float}
\usepackage{enumitem}
\usepackage{fancyhdr}
\usepackage{caption}
\usepackage{hyperref}

\definecolor{hyblue}{HTML}{1F4E79}
\definecolor{hyred}{HTML}{A23B3B}
\definecolor{hyorange}{HTML}{A66A00}
\definecolor{hygreen}{HTML}{2F6B4F}
\definecolor{hygray}{HTML}{666666}
\definecolor{hylight}{HTML}{EEF3F7}
\hypersetup{colorlinks=true,linkcolor=hyblue,urlcolor=hyblue}
\setlist[itemize]{leftmargin=1.7em,itemsep=0.2em}
\setlist[enumerate]{leftmargin=1.8em,itemsep=0.2em}
\renewcommand{\arraystretch}{1.22}
\captionsetup{font=small,labelfont=bf}
\emergencystretch=3em
\sloppy

\newcommand{\file}[1]{\texttt{#1}}
\newcommand{\sevA}{\textcolor{hyred}{\textbf{严重（A）}}}
\newcommand{\sevB}{\textcolor{hyorange}{\textbf{重要（B）}}}
\newcommand{\sevC}{\textcolor{hyblue}{\textbf{次要（C）}}}
\newcommand{\sevD}{\textcolor{hygray}{\textbf{观察（D）}}}
\newcommand{\fixed}{\textcolor{hygreen}{\textbf{已修复}}}
\newcommand{\clarified}{\textcolor{hyblue}{\textbf{契约澄清}}}

\pagestyle{fancy}
\setlength{\headheight}{14pt}
\fancyhf{}
\fancyhead[L]{src/model 数学模型审计报告}
\fancyhead[R]{HySim-XJTU-HRPES}
\fancyfoot[C]{\thepage}

\title{\textbf{模型层（\texttt{src/model} 与 \texttt{include/hacdcpf/model/}）数学模型\\工业级审计报告}\\[0.5em]
\large 最严苛工业应用视角 \textbar{} 全量逐行审计 \textbar{} 对照《元件模型技术手册》\textbar{} 数值实证}
\author{审计执行：AI 代理审计组（8 路并行审计 + 主审逐行复核）}
\date{2026 年 8 月 2 日}

\begin{document}
\maketitle

\begin{abstract}
\noindent
本报告对 HySim-XJTU-HRPES 平台模型层——\file{src/model}（4 个源文件、5\,336 行）与其组件契约头文件 \file{include/hacdcpf/model/}（约 4\,501 行）——的\textbf{数学模型}进行了最严苛视角的工业级审计。该层承担：组件 POD 语义、等值设备投影（变压器/开关/电机/聚合体）、rich$\to$canonical 投影（零阻抗合并、死岛剥离、能量路由器展开、编号规范化）、单位换算与参数库、以及结果归因（rich 索引空间回填）。

本文件是迁移前元件模型审计证据，原参考手册已不再作为当前规范源。关键公式曾经 Python/C++ 数值实证
（链接 \file{libhacdcpf.a} 端到端复现），并对严重级与关键重要级发现逐行核对源代码原文。本报告只收录
有代码行号与机理证据坐实的发现；当前行为以同目录 \file{model\_manual.tex} 与现行源码为准。

\textbf{结论与处置状态：}原始审计确认模型层核心等值/投影/换算主干正确，同时发现 \textbf{3 项严重（A 级）缺陷、25 项重要（B 级）缺陷与 20 项次要（C 级）缺陷}。本轮已对 A1--A3、B1--B25、C1--C20 逐项实施修复或收紧契约：不能由现有字段诚实表达的 Exponential 负荷改为显式拒绝；Charger 缺少母线、ER 端口域错误等情况改为 fail-closed；字段单位分歧统一到唯一消费口径。新增专用回归 \file{tests/test\_component\_models\_math\_audit.cpp}，并同步修正锁定旧错误语义的既有断言。修复后的核心库完整编译通过，专用回归 14 个用例、102 个断言通过；连同归因、参数、三相变压器、设备潮流、validation、ER 与短路回归，本轮共执行 10 个测试二进制、140 个用例、1\,574 个断言，全部通过。D 级项仍作为工程观察和后续维护债务保留，不在本轮宣称全部消除。
\end{abstract}

\tableofcontents
\newpage

\section{修复实施与验证记录}
\label{sec:remediation}

本节记录原始审计之后的代码处置。``已修复''表示模型行为或校验已改变；``契约澄清''表示既有计算口径保留，但结构体注释、参数规则或报告口径已统一。原始缺陷描述保留在后文，作为修复前证据。自动验证的主入口为
\file{tests/test\_component\_models\_math\_audit.cpp}；表中``静态核对''指对全部消费点做搜索和公式复核，不等同于端到端数值测试。

\subsection{A 级处置}

\begin{longtable}{@{}p{0.7cm}p{8.2cm}p{6.2cm}@{}}
\toprule
\textbf{编号} & \textbf{实施修改} & \textbf{验证}\\
\midrule
\endfirsthead
\toprule
\textbf{编号} & \textbf{实施修改} & \textbf{验证}\\
\midrule
\endhead
A1 & \fixed{}。\file{parity\_formulation.cpp} 在构建 Parity 问题时，从固定有功需求及其 ZIP 加权分子中扣除可控 FlexibleLoad 的 canonical 恒功率基线；无功基线保留，因为当前无 flexible-Q 决策。 & 专用回归构造 6~MW ZIP 负荷与 4~MW FlexibleLoad，断言固定 $P_d=6$~MW、$Q_d=4$~MVar，且 ZIP 比例恢复为 20/30/50；flex 变量仅保留一份。\\
A2 & \fixed{}。Charger 聚合改为按站点先清零再重算；缺失站点或站点母线无效抛 \file{std::invalid\_argument}；charger-only 因 Charger 本身无 bus 字段而明确拒绝，不再合成 \file{bus=0}。 & 专用回归以预填脏值站点和 60+40~kW 两台 Charger 验证结果为 100~kW/2 台及正确无功；charger-only 必须抛异常。\\
A3 & \fixed{}。ER 端口按 \file{port\_type} 分域：DC 端口直接复用引用的 DC 母线且不创建 VSC，只有 AC 端口创建 VSC；增加缺失母线、单侧端口、同侧多个 DC 母线、两侧同一 DC 母线守卫；溯源改为显式 expansion record。 & 专用回归以 AC/DC 同号可能冲突的域模型断言纯 DC 端口生成 0 台 VSC、DCDC 使用 canonical DC bus；另验证侧 B PQ 设定 $2-0.5=1.5$~MW 被原符号写入 \file{p\_ref\_mw}。\file{test\_multiport\_vpp\_er} 全部通过。\\
\bottomrule
\end{longtable}

\subsection{B 级处置}

\begin{footnotesize}
\begin{longtable}{@{}p{0.7cm}p{8.4cm}p{6.0cm}@{}}
\toprule
\textbf{编号} & \textbf{实施修改} & \textbf{验证}\\
\midrule
\endfirsthead
\toprule
\textbf{编号} & \textbf{实施修改} & \textbf{验证}\\
\midrule
\endhead
B1 & \fixed{}。2W 变压器投影统一执行 $Z_{1,0}/n_p$、$S_N n_p$，产物 \file{n\_parallel=1}，防止下游二次缩放；含 z0 缺省回退路径。 & 专用回归断言 $n_p=2$ 时正/零序 $x=0.5$~pu、容量 20~MVA。\\
B2 & \fixed{}。异步电机 \file{r\_pu/x\_pu} 统一定义为电机铭牌基准 pu；投影原值保存，四条短路路径按 $Z_{sys}=Z_{motor}S_b/S_N$ 换到系统基准；三份 IEC fixture 中误存的欧姆值按 $Z_{base}=V_N^2/S_N$ 迁移为 pu。 & 专用回归断言 motor-base pu 与容量不被改写；IEC 60909-4 示例恢复 $I_k''=19.52$~kA、其中电机贡献 4.78~kA，短路交叉回归通过。\\
B3 & \fixed{}。3W 三对阻抗按对应受端绕组 $V_N/V_{base}$ 平方折算，并把额定变比/母线基准失配写入每对 branch tap。 & 专用回归用 33/11 与 30/10~kV 两组铭牌，断言阻抗比 $(30/33)^2$、tap 为 1.1/1.1/1.0。\\
B4 & \fixed{}。三相变压器 fallback 投影复制 \file{tap\_side/tap\_min/tap\_max}。 & 专用回归断言 LV 侧与 $[-4,6]$ 档位边界不丢失。\\
B5 & \fixed{}。三相负荷电机以 $|Z''|=1/I_{LR}$、$R=X/(X/R)$ 重构，并携带总视在容量；去除 X/R 到电抗的量纲错配和容量二次缩放。 & 专用回归以 $I_{LR}=5$、$X/R=4$ 逐式比对 $R''$、$X''$ 与 0.9~MVA 容量。\\
B6 & \fixed{}。三相线路实际值换算统一使用投影后的 \file{out.ac.base\_mva}。 & 专用回归设置全局 100、三相 50~MVA，断言 10~kV、0.1+j0.2~$\Omega$ 投影为 0.05+j0.10~pu。\\
B7 & \fixed{}。删除半裕度猜测；DCDC Power 模式 \file{p\_ref} 改为侧 B AC-PQ 端口有符号设定值之和。 & A3 专用回归断言正负端口合计 1.5~MW。\\
B8 & \fixed{}。零阻抗合并增加 \file{tap=1} 与零移相守卫。 & 专用回归断言 $10^{-6}$~pu、tap=1.05 支路不合并。\\
B9 & \fixed{}。端子残差补齐 bus gs/bs、负荷清洗、renewable Q、电机、shunt、移动储能、VPP、微网及固定 VSC；同受端多个闭合设备等分不可辨识总残差。 & 专用回归以 4+j2 负荷和 1+j0.6 可再生电源验证两断路器各 1.5+j0.7；\file{test\_device\_flows} 通过。\\
B10 & \fixed{}。AsymmetricLoad 归因与投影共用非负有限 scaling。 & 专用归因回归断言 1~MW、scaling=2 报告 2~MW。\\
B11 & \fixed{}。归因最终统一把所有停运组件的 value 与终端 P/Q/V 置零。 & 专用回归断言停运 renewable 的 5~MW 署名值归因为 0；既有归因测试通过。\\
B12 & \fixed{}。PF 下同母线 switch/CB 无法识别物理切面时降为 AuditOnly，不再标 Strong。 & 归因恢复等级静态核对；既有 \file{test\_result\_attribution} 通过。\\
B13 & \fixed{}。设计手册补全执行 $R,X\div n_p$、rating$\times n_p$，三相同步写回沿用相同 pu。 & 专用参数回归比较 $n_p=1/2$，断言阻抗减半、容量加倍。\\
B14 & \fixed{}。相矩阵模式仅从矩阵计算阻抗；逐元素检查有限性与范围，NaN 不再被 \file{std::max} 掩盖。 & 专用回归注入矩阵 NaN，必须产生 \file{three\_phase\_line\_zero\_}\newline\file{series\_impedance}。\\
B15 & \clarified{}。\file{ACBranch::failure\_rate} 注释统一为整段元件的次/年；per-km 只存在于参数物化输入，物化后乘长度一次。 & 全库消费点静态搜索；可靠性参数回归通过。\\
B16 & \fixed{}。默认参数库与 TypicalParameters 的 DC 配电电压统一为 0.75~kV。 & 专用回归比较两条默认来源数值完全一致。\\
B17 & \fixed{}。VDC 模式仅在 $|k_{vdc}|>10^{-12}$ 或显式 grid-forming 时提供 DC 电压参考。 & 专用回归分别断言 $k=0$ 不提供、$k=1$ 提供参考。\\
B18 & \fixed{}。头文件契约改为与实现一致；DC static generator 与 AC 一样区分 controllable，固定零注入不再借用 nameplate 夸大容量。 & 专用回归断言同一设备不可控容量 0、可控容量 10~MW；Inf scaling fail-closed 为 0。\\
B19 & \fixed{}。Hydro 增加双向字符串分支。 & 专用回归断言 Hydro$\to$``Hydro''$\to$Hydro。\\
B20 & \fixed{}（兼容方案）。生命周期消费端把 $>1$ 的 \file{eol\_percent} 除以 100，保留 0.8 与 80 两种历史输入；结构体注释声明兼容规则。 & 公式静态核对并完整编译生命周期模块。\\
B21 & \fixed{}。\file{mag0\_percent} 与 \file{i0\_percent} 一致除以 100。 & \file{test\_three\_phase\_nr\_}\newline\file{transformer\_source} 24 用例、314 断言通过；对称相角断言按物理对称性而非旧错误幅值更新。\\
B22 & \clarified{}。保留代码的加性契约，DCBus 直挂需求与 DCLoad 均参与平衡；字段注释明确 additive。 & AC/DC 注入装配静态核对与 validation 回归。\\
B23 & \fixed{}（拒绝伪支持）。因 POD 无指数参数，validation 与 SolverData 对 Exponential 显式报错/抛异常，要求使用 ZIP 或 ConstantPower。 & 专用 validation 回归同时断言 \file{model} 诊断；核心装配编译通过。\\
B24 & \fixed{}。JSON 报告、市场、时序、可靠性与弹性统一 \file{rate\_a\_mva} 优先，\file{s\_max\_mva} 仅回退。 & 专用 JSON 回归令 rate=10、smax=5，10~MW 潮流必须报告 100\% 而非 200\%。\\
B25 & \fixed{}。字段统一声明为 \%/h；时序递推改为 $(1-p/100)^{\Delta t}$，UC 初值与 AC/DC storage 均使用同一函数。 & 公式静态核对、时序目标完整编译；1~h 时保持历史等价。\\
\bottomrule
\end{longtable}
\end{footnotesize}

\subsection{C 级处置}

\begin{footnotesize}
\begin{longtable}{@{}p{0.7cm}p{8.4cm}p{6.0cm}@{}}
\toprule
\textbf{编号} & \textbf{实施修改} & \textbf{验证}\\
\midrule
\endfirsthead
\toprule
\textbf{编号} & \textbf{实施修改} & \textbf{验证}\\
\midrule
\endhead
C1 & 储能仅在 $p_{mw}>0$ 放电时作为发电源；充电不再使死岛判活，投影与 island detector 同口径。 & 岛检测相关目标编译与既有回归。\\
C2 & SLACK/PV 死岛锚点必须 \file{in\_service=true}。 & validation/岛检测静态核对。\\
C3 & 合并后删除的自环支路把总 \file{b\_pu} 折为母线 \file{bs\_mvar}，不丢两端充电。 & 专用回归断言 0.2~pu 在 100~MVA 基准恢复为 20~MVar。\\
C4 & legacy 与四级 validation 均拒绝 bus index$\le0$、AC/DC 自环。 & 专用回归同时断言 \file{index} 与 \file{from\_bus/to\_bus} 错误。\\
C5 & attribution coverage 排除映射为 $-1$ 的已删支路。 & 专用 merge 回归断言已删两支路不再 total。\\
C6 & ER 溯源由名称前缀改为显式 \file{EnergyRouterExpansionRecord}。 & A3 投影映射静态核对。\\
C7 & VSC 短路字段与全部可靠性规则接入 \file{check\_range}。 & 专用回归以越界 \file{r\_sc/i\_max} 断言诊断存在。\\
C8 & VSC 短路阻抗规则描述明确为换流器自身额定功率/电压基准。 & 参数库规则静态核对。\\
C9 & 电压上下限冲突时只补缺失的一侧；两侧均由用户填写的冲突不被静默覆盖，留给 validation 报错。 & \file{test\_typical\_parameters} 通过。\\
C10 & \file{r\_ohm\_per\_km} 保存 20\textcelsius{} 手册值，pu 仍按运行温度热态电阻计算。 & 更新旧测试预期；\file{test\_typical\_parameters} 15 用例、401 断言通过。\\
C11 & 可靠性缺长度而采用 1~km 暴露时写入 \file{StandardParameterApplyReport::warnings}。 & 参数填充回归与源代码检查。\\
C12 & AC 阻抗任一分量缺失均单独补齐；缺失 \file{vkr} 回填后限制为不大于 \file{vk}。 & 参数库回归通过。\\
C13 & 仅提供 \file{c\_nf\_per\_km} 时按 $B=2\pi fC$ 生成 \file{b\_pu}；充电换算独立于串联阻抗是否已有 pu。 & 专用回归保留既有 r/x，并逐式断言 50~Hz 电容电纳。\\
C14 & 充电站容量改为 $n_fP_f+n_sP_s$，无分类台数时才回退总台数慢充。 & 专用回归 2 快+1 慢断言 127~kW。\\
C15 & 合法 \file{Load.scaling=0} 与 \file{Shunt.current\_step=0} 保留；只修负值、非有限值或越界值。 & 专用回归断言两字段填充后仍为 0。\\
C16 & DCStorage 与 AC storage 同样修复倒挂 SOC 窗口并重置越界初值。 & 专用回归输入 0.9/0.2，断言填充后窗口有序。\\
C17 & 未附着求解观测的行默认为 Unsupported，而非 AuditOnly；coverage 不再恒真。 & 专用与既有归因回归均断言 \file{total=false} 且 unsupported$>0$。\\
C18 & 多个同母线源断路器采用等分约定时降为 Approximate；单个仍可 Strong。 & 恢复等级静态核对。\\
C19 & AC/DC load、static generator、AsymmetricLoad 统一非负有限 scaling；负值、NaN、Inf 均按 0 fail-closed。 & 专用 helper/归因回归与 SolverData 编译。\\
C20 & validation 强制 ZIP 各项有限、非负、和为 100；装配前按每个负荷归一化，零和回退恒功率，防止绕过 validation 的内部调用改变 $P(1)$。 & 专用回归断言 30/30/30 被拒绝；A1 同时验证归一化后的 ZIP 混合。\\
\bottomrule
\end{longtable}
\end{footnotesize}

\subsection{构建口径与剩余边界}

本轮使用独立 Debug 构建目录 \file{/tmp/hacdcpf\_component\_audit\_build}，关闭 ETAP、Ipopt 与 SuiteSparse，以隔离可选依赖；因兄弟仓库 MIPSolvers 的本地提交/脏状态与记录 pin 不一致，显式使用非 Release 的 dirty-check override。该配置用于验证本项目修改，不等价于发布构建认证。专用测试不依赖 Ipopt；完整 \file{hacdcpf} 静态库已编译。最终实跑的 10 个测试二进制合计 140 个用例、1\,574 个断言，全部通过；其中 IEC 60909-4 与短路交叉回归分别为 8 与 502 个断言。D1--D20 仍是观察项，其中死 API、命名一致性、手册自动生成和跨模块默认值治理应进入后续维护，不应误读为本轮已全部解决。

\newpage

\section{审计范围、方法与定级标准}

\subsection{范围}

\file{src/model/network\_utils.cpp}（2\,746 行：等值投影、validation、零阻抗合并、死岛剥离、canonical 投影、ER 展开、端子潮流恢复）、\file{src/model/result\_attribution.cpp}（743 行：结果归因）、\file{src/model/standard\_parameter\_library.cpp}（1\,753 行：标准参数库与设计手册补全）、\file{src/model/unit\_conversion.cpp}（94 行：实际值$\to$标幺）、以及 \file{include/hacdcpf/model/} 全部头文件（组件 POD、defaults、effective\_capacity、device\_control\_role、typical\_parameters、enum 序列化）。

\subsection{方法}

\begin{enumerate}
  \item \textbf{逐行通读}：每份文件由独立审计通道全行通读（大文件分段双人），禁止抽样跳读。
  \item \textbf{手册对照}：以《元件模型技术手册》（pdftotext 全文可检索）为基准，对 41 种元件的等值电路、单位制、默认值、``已知边界''声明逐条核对实现；不一致处单列并标注手册章节。
  \item \textbf{数值实证}：约 15 个 C++ 探针（链接 \file{build/macos-release/libhacdcpf.a}，与 tests 同链）与多组 Python 复算；端到端复现（投影产物检查、\file{solve\_power\_flow} 全流程）用于坐实严重级发现。全部探针位于 \file{/tmp}，\textbf{未修改仓库任何文件}。
  \item \textbf{主审复核}：主审对全部 A 级与关键 B 级发现按行号重新阅读源代码原文，确认摘录、机理与定级无误后方予收录；涉及设计意图的标注``待核实''。
  \item \textbf{回归基线}：\file{test\_result\_attribution}、\file{test\_device\_flows}、\file{test\_typical\_parameters}、投影/合并系列（ctest \#48/\#303/\#575/\#771/\#1144/\#1145）审计时全部通过——第~\ref{sec:rootcause}~节解释为何既有测试未能暴露。
\end{enumerate}

\subsection{缺陷定级标准}

\begin{center}
\begin{tabular}{@{}llp{9.5cm}@{}}
\toprule
\textbf{级别} & \textbf{标记} & \textbf{定义}\\
\midrule
严重 & \sevA{} & 导致求解/研究结果物理错误且无告警，或投影层产出结构性错误模型。\\
重要 & \sevB{} & 确证的实现/语义错误但影响有界（特定数据形态、非默认路径触发），或手册/契约与实现系统性不符。\\
次要 & \sevC{} & 边缘情形、口径/文档不一致、校验缺口。\\
观察 & \sevD{} & 非数学错误，但应记录备查（工程取舍、潜在风险、文档漂移）。\\
\bottomrule
\end{tabular}
\end{center}

\section{严重缺陷（A 级）——必须严肃处理}
\label{sec:severe}

% ---------- A1 ----------
\subsection{A1　柔性负荷（FlexibleLoad）基准功率在 Parity OPF 有功平衡中被双重计数}
\label{sec:a1}

\begin{itemize}
  \item \textbf{位置}：投影 \file{src/model/network\_utils.cpp:106--125}（\file{:115--116} \texttt{ld.p\_mw = fl.p\_mw}）；清除清单 \file{:2432--2436}（\textbf{不含} \texttt{flexible\_loads}）；装配 \file{src/power\_flow/solver\_data.cpp:293,321}（\texttt{data.loads} 与 \texttt{data.flexible\_loads} 同时搬入）；消费 \file{src/optimal\_power\_flow/parity\_formulation.cpp:126--127}（\texttt{pd\_demand\_pu = data.pd\_pu}，含等值 Load）、\file{:792--793}（\texttt{p\_flex} 界 $[P_0-\Delta P_{dn},\,P_0+\Delta P_{up}]$）、\file{:1517--1519}（平衡方程）。
  \item \textbf{复核状态}：主审已逐行复核全部六处代码，机理链完整成立。
\end{itemize}

\textbf{机理。}投影把每个在运 FlexibleLoad 复制为恒功率 Load（$P_0$），rich 原件保留。OPF 有功平衡行
\begin{equation}
g_i = p_{calc} - p_{gen} - \cdots + {\color{hyred}p_d} - \Delta p_d + {\color{hyred}p_{flex,i}} - \cdots
\end{equation}
中，$p_d$ 经 \texttt{pd\_pu} 聚合\textbf{已含}等值 Load 的 $P_0$，而决策变量 $p_{flex}$ 以同一 $P_0$ 为界中心（暖启动即取 $P_0$）——初始点处需求被计为 $2P_0$（无功仅计一遍，P/Q 口径亦失配）。

\textbf{工业影响。}一切含柔性负荷的 AC OPF、需求响应与柔性平衡研究结果失真（基线需求翻倍）；潮流路径（PF 只消费 \texttt{loads}）不受影响；AsymmetricLoad 无此问题（SolverData 不携带原件）。

\textbf{实施结果。}采用 OPF 需求侧扣除等值 Load 的方案，同时重算扣除后的有功 ZIP 权重；FlexibleLoad 身份与变量继续保留，无功基线不重复建变量。验证见第~\ref{sec:remediation}~节 A1。

% ---------- A2 ----------
\subsection{A2　仅含充电机（Charger）建模时：站级功率双倍计数、且 \texttt{bus=0} 使全部充电功率静默丢失}
\label{sec:a2}

\begin{itemize}
  \item \textbf{位置}：合成路径 \file{src/model/network\_utils.cpp:2466--2489}（\texttt{cs.bus = 0} 于 \file{:2480}，预累加 \texttt{cs.p\_total\_kw += ch->p\_ch\_max\_kw} 于 \file{:2483}）；聚合 \file{:615--637}（\texttt{cs.p\_total\_kw += p\_kw} 于 \file{:630}、\texttt{cs.num\_chargers += 1} 于 \file{:635}）；丢弃点 \file{src/power\_flow/solver\_data.cpp:175--176}。
  \item \textbf{复核状态}：主审已逐行复核；两路独立审计通道分别给出实测（链接 \file{libhacdcpf.a}）：两台 60+40 kW 充电机 $\to$ \texttt{p\_total\_kw=200}、\texttt{num\_chargers=4}（应为 100/2）；120 kW 单台 $\to$ 240 kW。
\end{itemize}

\textbf{机理。}无 ChargingStation 时投影先自动建站并累加一遍 $\sum P_{ch}$（\file{:2483}），随后无条件调用 \file{project\_chargers\_into\_stations} 再对同一批充电机累加一遍 P（并首次累加 Q、台数再 $+1$/台）：
\begin{equation}
P = 2\sum P_{ch},\qquad Q = 1\sum P_{ch}\tan\varphi,\qquad N = 2N_{ch},
\end{equation}
P/Q 口径自相矛盾（隐含功率因数被篡改）。更致命的是 Charger 无 \texttt{bus} 字段，合成站点 \texttt{cs.bus = 0}，潮流装配 \texttt{idx = bus - 1 = -1} 直接跳过——该路径下\textbf{全部 EV 充电负荷静默从潮流中消失}，无任何警告。

\textbf{工业影响。}充电设施两条输入路径均出错：有站建模正常，无站建模（``只有充电桩清单''）负荷全丢；若日后修正 \texttt{bus=0}，双倍计数立刻显化为需求翻倍。

\textbf{实施结果。}聚合函数先重置站级 P/Q/台数再累加；无法由 Charger 确定母线时拒绝 charger-only 投影，缺站或无效站点母线同样抛异常。验证见第~\ref{sec:remediation}~节 A2。

% ---------- A3 ----------
\subsection{A3　能量路由器（ER）DC 端口被展开为 VSC，DC 域母线号写入 AC 域字段（域混淆）}
\label{sec:a3}

\begin{itemize}
  \item \textbf{位置}：端口划分 \file{src/model/network\_utils.cpp:1962--1970}（仅按 \texttt{port.side} 分侧，\textbf{不按} \texttt{port\_type} 分流）；VSC 创建 \file{:2020}（侧 A）与 \file{:2083--2086}（侧 B）：\texttt{vsc.bus\_ac = p->bus} 无条件执行。
  \item \textbf{复核状态}：主审已逐行复核；端到端实证（投影产物检查 $+$ \file{solve\_power\_flow}）：AC 母线 1(SLACK)/2(PQ) $+$ DC 母线 1，ER 端口 A=AC bus2、端口 B=DC bus1 $\to$ 投影产物 \texttt{VSC bus\_ac=1 bus\_dc=3}（DC 母线号 1 被当成 AC 母线 1，恰为平衡母线），求解``收敛''且 \texttt{p\_ac=+1.0} MW 注入平衡母线，用户 DC 母线成为无源孤岛，\texttt{er\_port\_transfers} 为空，无任何报错。
\end{itemize}

\textbf{机理。}\texttt{EnergyRouterPort::bus} 的域由 \texttt{port\_type} 决定（\file{converter\_components.hpp:265}）。对 DC 端口，\texttt{p->bus} 是 DC 母线编号，却被赋给 VSC 的 AC 域字段 \texttt{bus\_ac}；物理上 DC 端口应直连内部 DC 母线（或经独立 DCDC），根本不该生成 VSC。这直接违反项目铁律~1（AC/DC 域用 domain-qualified map，不混用）。若 DC 编号与 AC 编号不碰撞（如 99），装配守卫 \texttt{ac\_bus >= n} 会\textbf{静默丢弃}该端口功率而 DC 侧仍从 dc\_$\beta$ 抽取，ER 沦为纯吸能黑箱——两种情形均无报错。手册 \S 6.3（图 29 明示 DC 端口存在）``每端口一台 VSC''的表述照抄了该错误语义且未声明 DC 端口限制。

\textbf{工业影响。}含 DC 端口的能量路由器算例得到``收敛但物理错误''的潮流结果；错接功率可污染平衡母线功率平衡及全部下游分析（损耗、市场、可靠性）。

\textbf{实施结果。}DC 端口复用其 domain-qualified DC 母线，AC 端口才生成 VSC；非法跨域引用和不可构成两侧拓扑的 ER fail-closed。验证见第~\ref{sec:remediation}~节 A3。

\newpage
\section{重要缺陷（B 级）}
\label{sec:major}

\subsection{投影与等值层}

\subsubsection*{B1　\texttt{Transformer2W::n\_parallel} 在投影中被完全忽略}
\begin{itemize}
  \item \textbf{位置}：\file{network\_utils.cpp:155--235}（函数体零引用，已 grep 证实）；字段定义 \file{ac\_components.hpp:154}。\textbf{主审已复核。}
  \item \textbf{机理}：$n$ 台并联应阻抗 $\div n$、容量 $\times n$——平台自身在 \file{unit\_conversion.cpp:44--45} 对 \texttt{ACBranch.n\_parallel} 正是如此兑现，且 canonical 层 \texttt{ACBranch.n\_parallel} 有真实消费方。2W 投影生成的等值支路既不除也不乘，\texttt{n\_parallel} 保持默认 1。
  \item \textbf{影响}：并列变压器台数 $>1$ 时阻抗高估 $n$ 倍、热稳容量低估 $n$ 倍，潮流/短路/OPF 全面偏差；手册 \S 2.3 仅列字段未声明``仅元数据''。
\end{itemize}

\subsubsection*{B2　异步电动机 \texttt{r\_pu/x\_pu} 单位语义三方矛盾（代码=欧姆，手册+参数库=标幺）}
\begin{itemize}
  \item \textbf{位置}：\file{network\_utils.cpp:487--494}（\texttt{ld.r\_sc\_pu = m.r\_pu / z\_base\_motor}，注释 ``convert from ohms to pu''）。\textbf{主审已复核。}
  \item \textbf{机理}：手册 \S 3.6 声明 \texttt{r\_pu/x\_pu} 为额定基准标幺（典型 0.01--0.15/0.1--0.25），\file{typical\_parameters.hpp:832--835} 按 pu 量级填充（0.002/0.2）；而投影与短路专用路径（\file{short\_circuit.cpp:525--529}）均按\textbf{欧姆}消费（除以 $Z_{base}=V_n^2/S_n$）。IEC 60909-4 验证算例恰按欧姆填写，故验证通过而语义矛盾未暴露。
  \item \textbf{影响}：用户按手册填 0.025+j0.25 pu 时，阻抗被再除 $Z_{base}$（典型 6 倍），电机馈供短路电流同倍\textbf{高估}。修复需项目层面裁决单位制后统一改手册或代码。
\end{itemize}

\subsubsection*{B3　\texttt{Transformer3W} 投影缺电压基准折算（2W 有、3W 无）}
\begin{itemize}
  \item \textbf{位置}：\file{network\_utils.cpp:270--298}（仅 \texttt{rx\_from\_vk\_vkr(..., base\_mva, sn\_pair)} 容量折算）；对照 2W 的 \file{:177--187}（$(V_n/V_{base})^2$）与 \file{:199--207}（额定变比失配项）。\textbf{主审已复核。}
  \item \textbf{影响}：母线 \texttt{base\_kv} $\neq$ 绕组 \texttt{vn\_kv}（MATPOWER 导入常见）时，三绕组阻抗错 $(V_n/V_{base})^2$ 倍、变比缺失配项；校验层不查电压一致性，手册 \S 2.4 未声明``$V_n$ 须等于母线基准''。
\end{itemize}

\subsubsection*{B4　三相变压器 $\to$ Transformer2W 投影丢 \texttt{tap\_side}（及 \texttt{tap\_min/tap\_max}）}
\begin{itemize}
  \item \textbf{位置}：\file{network\_utils.cpp:698--721}（复制 \texttt{tap\_pos}、\texttt{tap\_neutral}、\texttt{tap\_step\_percent}、\texttt{shift\_deg}、\texttt{z0}，无 \texttt{tap\_side}）。\textbf{主审已复核。}
  \item \textbf{影响}：\texttt{tap\_side=1}（低压侧分接头）的三相变压器按高压侧处理，变比方向（$t$ vs $1/t$）与阻抗 $t^2$ 归算同时错误。
\end{itemize}

\subsubsection*{B5　\texttt{ThreePhaseLoad}$\to$Load 电动机馈供字段错配（\texttt{x\_sub\_pu} $\leftarrow$ \texttt{x\_r\_ratio}）}
\begin{itemize}
  \item \textbf{位置}：\file{network\_utils.cpp:739--740}。\textbf{主审已复核。}
  \item \textbf{机理}：\texttt{Load.x\_sub\_pu} 是次暂态电抗 $X''$（pu，短路馈供用），\texttt{ThreePhaseLoad.x\_r\_ratio} 是无量纲 $X/R$ 比（典型 3--10）——量纲错配；\texttt{lrc\_pu}（$|Z|\approx 1/I_{LR}$ 的唯一来源）未折算，\texttt{eq.sn\_mva} 未赋值（默认 0）$\to$ 短路装配 \texttt{sn\_mva<=1e-6} 直接跳过，三相负荷电机馈供\textbf{静默缺失}；一旦补上 \texttt{sn\_mva}，又以 10--50 倍过大的``电抗''参与计算。正确形式：$X'' = 1/\texttt{lrc\_pu}$、$R = X''/\texttt{x\_r\_ratio}$，并携带容量基准。
\end{itemize}

\subsubsection*{B6　\texttt{project\_three\_phase\_if\_needed} 基准功率不自洽}
\begin{itemize}
  \item \textbf{位置}：\file{network\_utils.cpp:644}（\texttt{out.ac.base\_mva = tp.base\_mva ?: out.base\_mva}）vs \file{:681}（$Z_{base}=kV^2/\texttt{out.base\_mva}$）。
  \item \textbf{影响}：下游（短路/DC OPF/谐波/JSON）优先用 \texttt{ac.base\_mva}，而线路欧姆$\to$标幺用全局 \texttt{base\_mva}；两者不等时三相系统阻抗基准整体错位。默认两者均 100，属潜伏缺陷。
\end{itemize}

\subsubsection*{B7　ER 展开 DCDC Power 模式 \texttt{p\_ref} 启发式任意猜测，直接决定侧 B 构网/下垂端口运行点}
\begin{itemize}
  \item \textbf{位置}：\file{network\_utils.cpp:2151--2163}（\texttt{remaining * 0.5} 裕度猜测）。
  \item \textbf{机理}：DCDC 是 dc\_$\alpha\leftrightarrow$dc\_$\beta$ 唯一功率通道，Power 模式 \texttt{p\_ref} 固定后，侧 B 的 VF/Droop VSC 只能被动取 $p_{ref}-\sum p_{PQ}$——0.5 裕度猜测直接设定该端口物理传输功率；负 \texttt{p\_set}（受电）端口被排除在求和外，系统性偏置。手册 \S 6.3 未声明此启发式；现有测试只断言收敛从未断言 droop 端口功率。
\end{itemize}

\subsection{拓扑合并（潜伏）}

\subsubsection*{B8　公共零阻抗合并入口对含分接头/移相支路无任何防护（潜伏缺陷）}
\begin{itemize}
  \item \textbf{位置}：\file{network\_utils.cpp:1590--1593}（判据只查 $|r|,|x|,|b|$）；公共入口 \file{:1831--1834}。\textbf{主审已复核判据代码。}
  \item \textbf{机理}：带非额定变比/相移的支路两端\textbf{不是等电位点}，并入同一等价类等于把理想变压器/移相器短接。实证：\texttt{tap=1.05}、$r=x=10^{-6}$ 支路经公共路径后两母线被收缩为 1 个。
  \item \textbf{暴露面校准}：投影主流水线走白名单（仅开关/断路器展开支路，恒为 tap=1/shift=0），\textbf{生产路径安全}；公共 API \file{merge\_zero\_impedance\_buses} 当前全仓库无生产调用点（已 grep 证实），属潜伏风险——任何外部调用方/后续模块用该 API 处理含低阻抗变压器的系统即静默得到错误拓扑。修复：补 \texttt{tap != 1.0 || shift\_deg != 0.0} 防护。
\end{itemize}

\subsection{结果归因层}

\subsubsection*{B9　\texttt{compute\_device\_terminal\_flows} 残差口径三项缺陷，``exactly'' 声明不实（归因层继承并标 Strong）}
\begin{itemize}
  \item \textbf{位置}：\file{network\_utils.cpp:2505--2515}（多开关）、\file{:2545--2603}（注入枚举）、\file{:2570--2574}（署名值）；继承方 \file{result\_attribution.cpp:583--606}。\textbf{主审已复核 (a)(b) 两处。}
  \item \textbf{机理}：(a) 同一受端母线 $N$ 台闭合开关\textbf{各自分得全额残差}（实证：两台开关各报 $-40$ MW，合计 $-80\neq-40$），KCL 在归因层被破坏；(b) 注入枚举不含 motors/shunts/母线 gs/bs/external\_grids/VSC/LCC/VPP/微网/ER 端口，且 \textbf{renewable 只加 P 漏 Q}（\file{:2580--2582}）——实证 renewable 发 $1.0+j0.5$ 经开关送出，归因 $p=1.0$（对）$q=0$（应 0.5）；(c) 发电机按署名 \texttt{pg\_mw} 而非求解输出计入，slack 拾取量与 PV 无功全部残差化到开关潮流。
  \item \textbf{影响}：变电站母联/分段开关潮流与 \texttt{loading\_pct} 在含电机、并联补偿、外网、换流器的典型工业拓扑下系统性失真，可误导 N-1 与重构决策；注释自称 ``\textbf{exactly} the collapsed device's terminal injection'' 不成立。修复：枚举完备化（或用求解注入）、多设备共母线声明``仅总和可恢复''并降级 Approximate。
\end{itemize}

\subsubsection*{B10　\texttt{asymmetric\_load} 归因漏乘 \texttt{scaling}，与投影端差一个因子}
\begin{itemize}
  \item \textbf{位置}：\file{result\_attribution.cpp:409--410}（不乘）；对照投影 \file{network\_utils.cpp:136--138}（乘）与开关残差法 \file{:2561}（乘）。\textbf{主审已复核两处。}
  \item \textbf{实证}：scaling=2.0、三相和 2.0 MW $\to$ canonical 等值负荷 $p=4.0$，归因报 $p=2.0$。下游功率平衡诊断（\file{run\_gui\_server.cpp:282--288}）正是消费此值，制造虚假平衡残差。
\end{itemize}

\subsubsection*{B11　\texttt{in\_service} 口径不一致：13 类退出运行元件归因非零铭牌/经验值}
\begin{itemize}
  \item \textbf{位置}：\file{result\_attribution.cpp}——检查了 \texttt{in\_service} 的仅 6 处（load/dc load/ac/dc static gen/dc pv）；未检查的 13 类含 storage、renewable\_gens（\file{:387--391}）、pv\_system、flexible\_load、asymmetric\_load、motor、shunt、charging\_station、mobile\_storage、vpp、microgrid、dc storage 两族。\textbf{主审已抽查复核。}
  \item \textbf{实证}：停运 renewable\_gen 归因 $p=1.0$；停运 pv\_system 按经验曲线报 $1.477$ MW，而求解端注入为 0。库内平衡诊断靠 \texttt{row.in\_service} 跳过故无恙，其它消费者无此保护。
\end{itemize}

\subsubsection*{B12　PF-only 路径下同母线源侧断路器静默归因 0 潮流且标 Strong}
\begin{itemize}
  \item \textbf{位置}：\file{result\_attribution.cpp:580--607}（设备流块）$+$ \file{:613--654}（修正启发式要求 \texttt{opf\_result} 非空，PF 生产调用点均传 \texttt{nullptr}）。
  \item \textbf{实证}：物理载 5 MW 的同母线断路器，归因 from/to $p=0$、recovery=\textbf{Strong}。GUI 潮流快照显示 0 MW/0\% 负载率，误导运行人员；按铁律应降级 AuditOnly 或注入诊断。
\end{itemize}

\subsection{参数库与设计手册补全}

\subsubsection*{B13　设计手册补全忽略 \texttt{n\_parallel}：阻抗放大 $n_p$ 倍、容量缩小 $n_p$ 倍}
\begin{itemize}
  \item \textbf{位置}：\file{standard\_parameter\_library.cpp:1229--1233}。\textbf{主审已复核。}
  \item \textbf{机理}：\texttt{new\_r\_pu = resistance * length\_km / zbase} 未除 \texttt{n\_parallel}——同平台权威折算 \file{unit\_conversion.cpp:44--48} 明确 \texttt{/ n}；手册公式同为 $R'l/(n_p Z_{base})$。实证：$n_p=2$ 时补全 $r_{pu}=0.1282$ vs 正确 $0.0641$（恰 2 倍）。ETAP/BPA 导入模型带 \texttt{n\_parallel}（\file{etap\_io.cpp:831}、\file{bpa\_io.cpp:547}），双回线路经补全后压降/网损偏大 2 倍、OPF 容量偏紧 2 倍。三相同步段（\file{:1333--1341}）连带污染 $r_0/x_0$。
\end{itemize}

\subsubsection*{B14　三相相矩阵线路的零阻抗校验存在掩盖漏洞（\texttt{std::max} 吞 NaN），且相矩阵元素从不做范围检查}
\begin{itemize}
  \item \textbf{位置}：\file{standard\_parameter\_library.cpp:822--841}。\textbf{主审已复核。}
  \item \textbf{机理}：(a) 校验量取 $\max(\texttt{hypot}(r_1,x_1),\ \texttt{hypot}(\text{各矩阵元}))$——\texttt{use\_phase\_matrix=true} 时求解器消费的是 9 元素相矩阵，与 $r_1/x_1$ 取 max 无物理意义；C++ 语义 \texttt{std::max(a, NaN)=a}，全零或含 NaN 的相矩阵被静默掩盖（已实测确认）；(b) \texttt{else if (!use\_phase\_matrix)} 使矩阵元素\textbf{从不进入} \texttt{check\_range}。误配置模型得到虚假``参数校验通过''，奇异相域导纳阵直到三相求解器才暴露。
\end{itemize}

\subsubsection*{B15　\texttt{ACBranch::failure\_rate} 单位语义手册与代码相反（per-km vs 绝对值/年）}
\begin{itemize}
  \item \textbf{位置}：字段 \file{ac\_components.hpp:95}；手册 \S 2.1 声明``次/(km·a)（乘长度）''；消费端 \file{reliability\_assessment.cpp:184}（直接作 \texttt{failure\_rate\_per\_year}）、\file{three\_stage\_reliability.cpp:324} 等\textbf{均不再乘长度}；仅填充路径 \file{standard\_parameter\_library.cpp:1002--1006} 做 $\lambda'\times l$。
  \item \textbf{影响}：按手册手填 per-km 值的用户，可靠性指标（SAIFI/SAIDI/EENS）被欠计线路长度倍（如 5 km 线路 5 倍），无任何诊断。修正对象主要是手册与字段注释（代码语义自洽：字段存绝对值/年）。
\end{itemize}

\subsubsection*{B16　\texttt{dc\_bus.base\_kv} 默认值在两条官方填充路径间相差 2 倍且未声明}
\begin{itemize}
  \item \textbf{位置}：\file{standard\_parameter\_library.cpp:409--412}（distribution 剖面 1.5 kV）vs \file{typical\_parameters.hpp:27}（一律 \texttt{kDCDistributionKv = 0.75}）。\textbf{主审已复核两处。}
  \item \textbf{影响}：同一稀疏模型经两条``标准''路径得到 $2\times$ 直流电压基准，DC 支路 $Z_{base}=V^2/S$ 随之相差 $4\times$，DC 潮流/短路结果不可比；手册 \S 1.3 把两者同列为``典型范围来源''未声明分歧。
\end{itemize}

\subsection{头文件契约}

\subsubsection*{B17　\texttt{provides\_dc\_v\_reference} 对下垂模式无条件置真，与同文件 \texttt{forms\_dc\_voltage} 的 \texttt{k\_vdc} 门控直接矛盾}
\begin{itemize}
  \item \textbf{位置}：\file{device\_control\_role.hpp:81--99}（VDC\_Q/VDC\_VAC/DC\_V\_DROOP\_AC\_V 无条件置真）vs \file{:62--66}（\texttt{forms\_dc\_voltage} 要求 \texttt{k\_vdc != 0}）。\textbf{主审已复核。}
  \item \textbf{影响}：\texttt{k\_vdc=0} 的 VDC\_Q 换流器：角色输出 \texttt{provides\_dc\_v\_reference=true} 且 \texttt{is\_dc\_grid\_forming=false}，同一函数自相矛盾；AC/DC OPF 装配（\file{acdcopf\_builder.cpp:118,161}）据其置 \texttt{is\_vdc\_slack}，而潮流协调校验（\file{converter\_coordination.cpp:385--399}）对同一换流器报 Error——两引擎对同一设备可行域结论相反。
\end{itemize}

\subsubsection*{B18　\texttt{effective\_capacity.hpp} 文件头契约与实现两处矛盾；\texttt{StaticGeneratorDC} 版忽略 \texttt{controllable}}
\begin{itemize}
  \item \textbf{位置}：约定块 \file{:28--32}（``no scaling''、``$|effective\_p\_mw|$''）vs 实现 \file{:100,:103--105}（乘 \texttt{scaling}、用 \texttt{max(0,·)}）；DC 版 \file{:116--122}（无 \texttt{controllable} 分支，AC 版 \file{:94--106} 有）。
  \item \textbf{影响}：该文件存在意义即``single source of truth''，契约与实现矛盾使 scaling$<$1 算例中可靠性/弹性评估的可恢复容量与 OPF 调度上限系统性不一致；DC 侧对 \texttt{controllable=false} 且 $p_{set}\approx0$、$p_{max}>0$ 的电源高估可恢复容量。
\end{itemize}

\subsection{组件 POD 语义（手册对照）}

\subsubsection*{B19　\texttt{RenewableType::Hydro} 序列化往返损坏（Hydro $\to$ ``Wind''）}
\begin{itemize}
  \item \textbf{位置}：\file{enum\_strings.hpp:219--230}（无 Hydro 分支，default 落 ``Wind''）；枚举声明 \file{enums/renewable\_enums.hpp:10} \texttt{Hydro = 3}。\textbf{主审已复核两处。}
  \item \textbf{影响}：\file{case\_builders.cpp:3014} 的内置算例真实使用 Hydro，JSON 导出被静默改写为 Wind，时序/碳核算分类统计被篡改（编译验证：\texttt{Hydro -> "Wind" -> Wind}，FAIL）。
\end{itemize}

\subsubsection*{B20　\texttt{eol\_percent} 字段名（\% 后缀）与消费语义（小数）矛盾}
\begin{itemize}
  \item \textbf{位置}：\file{ac\_components.hpp:599}（\texttt{double eol\_percent\{0.8\}}，Storage/MobileStorage/DCStorage 三处）；消费方 \file{lifecycle\_simulation.cpp:476}（\texttt{soh <= eol\_percent}，SOH 为 0--1 小数）。\textbf{主审已复核。}
  \item \textbf{影响}：手册 \S 3.5 自证矛盾（单位列 ``\%''、范围 ``60--90''、默认 0.8）；用户按字段名填 80 时 \texttt{soh $\le$ 80} 恒真，储能每年立即更换，寿命仿真结论被放大 100 倍。修复：更名 \texttt{eol\_fraction} 或消费方 \texttt{/100}。
\end{itemize}

\subsubsection*{B21　\texttt{mag0\_percent} 未除 100，与同结构体 \texttt{i0\_percent} 口径相反}
\begin{itemize}
  \item \textbf{位置}：消费方 \file{three\_phase\_nr.cpp:1275}（\texttt{mag0\_abs = |z0| * mag0\_percent} 直乘）vs \file{:1189}（\texttt{i0\_percent / 100.0}）。\textbf{主审已复核两处。}
  \item \textbf{影响}：同一结构体两个 \texttt{\_percent} 字段两种口径；手册 \S 8.3 以``已知边界''声明该怪癖但 JSON/schema 层无保护——按字段名填 2（\%）得到 200\% 零序励磁阻抗，接地故障研究可差 100 倍。
\end{itemize}

\subsubsection*{B22　\texttt{DCBus.pd\_mw} 手册声明与代码直接矛盾（``忽略'' vs 加性）}
\begin{itemize}
  \item \textbf{位置}：代码 \file{pf\_injection\_assembly.cpp:102--116}（注释 ``DC bus demand and explicit DCLoad records are additive''，无条件叠加）；手册 \S 5.1 声明``负荷表非空时 \texttt{pd\_mw} 被忽略（不与 DCLoad 叠加）''。
  \item \textbf{影响}：二者必有一错。AC 侧经 H1 已是加性，代码注释表明 DC 加性为有意设计，手册按过时行为撰写；按手册建模（母线直挂 $+$ DCLoad 并存）的用户将被\textbf{重复计量直流负荷}。
\end{itemize}

\subsubsection*{B23　\texttt{LoadModel::Exponential} 为不可实现的死枚举值}
\begin{itemize}
  \item \textbf{位置}：\file{enums/load\_enums.hpp}；\texttt{Load} 无指数 $\alpha/\beta$ 字段；稳态电压依从完全由 \texttt{z/i/p\_percent} 驱动（\file{solver\_data.cpp:160--168}），不读 \texttt{model} 枚举；手册 \S 3.1 仍将 Exponential 列为可用模型。
  \item \textbf{影响}：用户设 \texttt{model=Exponential} 得到的是恒功率（默认 ZIP 系数），无任何报错——电压依从研究可静默得出恒功率结论。
\end{itemize}

\subsubsection*{B24　\texttt{DCBranch.rate\_a\_mva} 与 \texttt{s\_max\_mva} 双额定字段跨模块优先级相反}
\begin{itemize}
  \item \textbf{位置}：\file{json\_io.cpp:48--53}（\texttt{s\_max} 优先）vs \file{market\_simulation.cpp:734--739}、\file{time\_series\_pf.cpp:894--895}、\file{three\_stage\_reliability.cpp:1775--1776} 等（均 \texttt{rate\_a} 优先）。
  \item \textbf{影响}：两字段同义（手册 \S 5.2）且未声明优先级；同时填写且不等时，限值执行模块用 \texttt{rate\_a} 而报告 \texttt{loading\_pct} 用 \texttt{s\_max}——``执行未越限、报告越限''或反之。
\end{itemize}

\subsubsection*{B25　\texttt{self\_discharge\_pct} 时间基准跨模块歧义（每时段 vs 每小时）}
\begin{itemize}
  \item \textbf{位置}：\file{time\_series\_pf.cpp:2312--2313}（\texttt{retention = 1 - pct/100}，每时段一次、与 $\Delta t$ 无关）vs \file{BasicDynamicDevices.cpp:9000--9001}（\texttt{/3600}，即 \%/h 连续）；\file{case\_builders.cpp:4352} 注释 ``0.01\%/hr''；手册 \S 3.5/\S 5.4（``每时段''）。\textbf{主审已复核时序侧。}
  \item \textbf{影响}：同一字段三种口径并存；UC 步长 $\Delta t\neq 1$ h（如 15 min）时，时序层比暂态层衰减快 4 倍；$\Delta t=1$ h 时两者一致，故缺陷长期隐蔽。修复：字段或文档声明时间基准（建议 \%/h），时序层按 $(1-\texttt{pct}/100)^{\Delta t}$ 折算。
\end{itemize}

\newpage
\section{次要缺陷（C 级）汇总}
\label{sec:minor}

\begin{footnotesize}
\begin{longtable}{@{}p{0.7cm}p{5.2cm}p{9.3cm}@{}}
\toprule
\textbf{\#} & \textbf{位置} & \textbf{问题}\\
\midrule
\endfirsthead
\toprule
\textbf{\#} & \textbf{位置} & \textbf{问题}\\
\midrule
\endhead
C1 & \file{network\_utils.cpp:1304--1307} & 死岛判据 \texttt{|p\_mw|>1e-9} 把\textbf{纯充电储能}误判为电源（手册 \S 3.5：$p<0$ 为充电负荷）；潮流侧 \file{island\_detector.cpp:270} 为同一判据，两层口径一致地错。\\
C2 & \file{network\_utils.cpp:1279--1283} & 死岛判据的 SLACK/PV 母线检查不验证 \texttt{in\_service}——停运 SLACK 母线使整岛判活，其负荷被装配静默置零且不在 \texttt{dead\_bus\_indices} 中，无诚实声明。\\
C3 & \file{network\_utils.cpp:1755--1759} & 自环支路（真实线路被并联零阻抗路径短接）删除时丢失两端 $B/2$ 充电功率，未折算为母线并联。\\
C4 & \file{network\_utils.cpp:935--1068} & \texttt{validate()} 放过 \texttt{bus.index}$\le 0$（下游按 \texttt{index-1} 寻址将越界）与自环支路；该版本无生产调用点（生产走四级 validation），但现代 validation 同样不查自环与 index$\le 0$。\\
C5 & \file{network\_utils.cpp:2653--2659} & \texttt{evaluate\_attribution} 未排除 \texttt{branch\_orig\_to\_proj=-1}（已删除分支），误报可恢复/Strong（实证：被删线路仍计 attributed）。\\
C6 & \file{network\_utils.cpp:2308--2320} & ER$\to$VSC/DCDC 溯源按名称前缀匹配，空名 ER（前缀 "\_"）或前缀包含关系导致交叉归因（仅溯源层面）。\\
C7 & \file{standard\_parameter\_library.cpp:}\newline\file{581--866} & 校验规则覆盖缺口：\texttt{vsc.r\_sc/x\_sc/i\_max} 与全部 28 条可靠性规则定义后从不参与 \texttt{check\_range}，越界值无诊断。\\
C8 & \file{standard\_parameter\_library.cpp:}\newline\file{423--428} & \texttt{vsc.r\_sc\_pu/x\_sc\_pu} 规则描述未声明标幺基准（消费端按换流器自身额定换算，\file{short\_circuit.cpp:604}）；按系统基准填写的用户短路贡献错 $S_{rated}/S_{base}$ 倍。\\
C9 & \file{standard\_parameter\_library.cpp:}\newline\file{894--899} & \texttt{vmin$\ge$vmax} 修复无条件覆盖两个字段（超出``仅补缺''语义），authored 信息丢失。\\
C10 & \file{standard\_parameter\_library.cpp:}\newline\file{1303} & 补全写回的 \texttt{r\_ohm\_per\_km} 为运行温度热态值（$+19.7\%$--$28.2\%$），与手册``20 \textcelsius{} 导线手册值''字段语义冲突（数值自洽，下游对照 GB/T 3956 时高估）。\\
C11 & \file{standard\_parameter\_library.cpp:}\newline\file{1002--1006} & 可靠性长度暴露缺省 1 km 为隐式假设，未在报告中声明（短线路系统故障率系统性偏高）。\\
C12 & \file{standard\_parameter\_library.cpp:}\newline\file{908--911,:924--933} & AC 支路阻抗``双缺才填''，单侧缺失留下 $x=0$ 空洞（DC 潮流 $B'$ 静默丢支路）；\texttt{vkr} 缺省回填可产生 \texttt{vkr>vk} 的物理不可能组合。\\
C13 & \file{unit\_conversion.cpp:52--56} & \texttt{c\_nf\_per\_km} 在 C++ 模型层从不参与标幺换算（仅 JSON 读写），手册声明的换算路径不实——只给电容参数的算例静默丢失全部充电电纳。\\
C14 & \file{typical\_parameters.hpp:801--804} & 充电站 \texttt{max\_power\_kw} 占位填充 $=$ 台数 $\times$ 慢充功率，忽略快充（2 台快充：14 kW vs 应 120 kW）。\\
C15 & \file{typical\_parameters.hpp:842,}\newline\file{:628,:647} & 占位填充改写合法工程状态：shunt \texttt{current\_step=0}（全部切出）改投 1 档；负荷 \texttt{scaling=0}（停运）重置为 1 并重新注入——与 \file{effective\_capacity.hpp} 的``scaling$\le$0 按 0''契约直接冲突。\\
C16 & \file{typical\_parameters.hpp:538--562} & DCStorage 填充路径缺 SOC 窗口倒挂修复（AC 侧 \file{:510--514} 有），倒挂输入窗口仍为空集。\\
C17 & \file{result\_attribution.cpp:}\newline\file{93,:712--729} & coverage 指标退化：\texttt{add\_collection} 一律 AuditOnly、全文件无 Unsupported 赋值点 $\to$ \texttt{coverage.total()} 恒真（测试 \file{test\_result\_attribution.cpp:149} 把退化行为锁定）。\\
C18 & \file{result\_attribution.cpp:641--652} & 同母线多断路器等分外部电网功率后标 Strong（理想化后电流分配本不确定，应标 Approximate 或声明约定）。\\
C19 & \file{solver\_data.cpp:34,:151} vs\newline\file{pf\_injection\_assembly.cpp:}\newline\file{114,:132} & 负 \texttt{scaling} 清洗口径跨域不一致：AC 负荷/电源裸乘（反向注入）、DC 负荷置零、DC 电源裸乘、AsymmetricLoad 投影第三口径（$\le 0 \to 1.0$）。\\
C20 & \file{solver\_data.cpp:160--168} & ZIP 系数之和不归一化、无校验（30/30/30 $\to$ $P(1.0)=0.9P_0$ 静默改变负荷量；手册声明``和应为 100''但代码不校验）。\\
\bottomrule
\end{longtable}
\end{footnotesize}

\section{观察项（D 级）汇总}
\label{sec:obs}

\begin{itemize}
  \item \textbf{D1}　电机零序参数 \texttt{r0\_pu/x0\_pu} 投影丢失（\file{network\_utils.cpp:454--497} 未映射；接地故障下电机零序馈供消失，多数电机中性点不接地故实际影响有限）。
  \item \textbf{D2}　三绕组星形腿近零 fallback 丢弃相角（\file{:290--298}，$|z|\le10^{-12}$ 时取纯实数；极限情形数值验证仍收敛正确值）。
  \item \textbf{D3}　\file{rx\_from\_vk\_vkr} 对 \texttt{vkr>vk}/负值非法数据静默钳 \texttt{x=0}（\file{:78--81}）；\texttt{bus\_base\_kv\_or\_default} 对缺失 \texttt{base\_kv} 回退 1.0 kV 不告警（\file{:32--39}）。
  \item \textbf{D4}　2W 未给 \texttt{z0\_percent} 时回退 $Z_0=Z_1$（\file{:226--229}）——手册 \S 2.4 对三绕组声明该回退、\S 2.3 对双绕组未声明，声明缺口。
  \item \textbf{D5}　合并代表母线选择不偏好 in-service 母线；成员母线的电压限值/\texttt{base\_kv}/初值无声丢失（\file{:1670--1701}，无最紧约束归并）。
  \item \textbf{D6}　VSC``有源''判据过宽（\file{:1308--1311}：任何在役 VSC 挂接即电源，不看控制模式与 DC 侧支撑）；\file{strip\_dead\_islands} 头注释与实现判据措辞不一致（\file{:1209--1220} vs \file{:1286,:1306}）。
  \item \textbf{D7}　幂等重投影在非 canonical 母线编号下会覆盖既有 BusMergeMap provenance（\file{:1626--1828}，\textbf{待核实}可达性；canonical 编号下早退无此问题）。
  \item \textbf{D8}　ER 展开：VF 端口 PV 提升后 \texttt{q\_set\_mvar} 静默失效；\texttt{v\_set\_pu} 一参三用（DC 设定/AC PV/换流器 AC 设定）；侧 A 自动升格与头注释不一致，[PQ, Droop] 端口序产生双定压源；单侧 ER 产生悬空内部 DC 母线；\texttt{num\_ports} 与 \texttt{ports.size()} 无校验（\file{:1947--2163}）。
  \item \textbf{D9}　\texttt{AttributedTerminal} 的 $p$ 符号双轨制（网络类``母线$\to$器件'' vs 源荷类``注入母线''）未在头文件声明；AC 母线 \texttt{recovery\_reason} 写 ``BusMergeMap voltage broadcast'' 但实际从不使用 BusMergeMap，占位值静默回退（\file{result\_attribution.cpp:238,:107--121}）。
  \item \textbf{D10}　DC 支路归因阈值（\texttt{r\_pu > 1e-12}）严于求解器（\texttt{r\_pu == 0} 跳过）；MATPOWER 提取的 2W 变压器在 ac\_branch 与 transformer\_2w 双重列报。
  \item \textbf{D11}　手册文档漂移：\S 6.3 对 \file{network\_utils.cpp} 的行号引用全部过期（约 72 行）；``48 条工程校验规则''实为 52 条；\S 2.1 默认值列（r\_pu 0.01/x\_pu 0.04、mttr 架空 8/电缆 24）与结构体初值 0.0 冲突；\S 5.1 \texttt{DCBus.pd\_mw} 按过时行为撰写（见 B22）。
  \item \textbf{D12}　\texttt{defaults.hpp:56} \texttt{kPackageVersion="0.5.0"} 与 \file{CMakeLists.txt:3} \texttt{VERSION 0.1.0} 两套版本出处；\texttt{kGenPmax/kGenPmin/kGenQmax/kGenQmin} 为死常量（全库无消费方，``单一来源''声明失实）。
  \item \textbf{D13}　\texttt{element\_type} 同名跨文件不同型（Switch 为 \texttt{int}、CircuitBreaker/DCCircuitBreaker 为 \texttt{std::string}）；\texttt{StaticGenerator::co2\_emission\_rate} 与 \texttt{emission\_factor\_tco2\_mwh} 同义异名；\texttt{VSCConverter::mtbf\_hr} 注释悬空引用不存在的 \texttt{Generator::mtbf\_hr}。
  \item \textbf{D14}　可靠性默认值在参数库与 TypicalParameters 间差异大（1 km 线路故障率相差 6 倍）；AC 支路 per-km$\times$长度 vs DC 支路固定 MTBF 的暴露建模不一致；架空线 0.50 occ/(km·yr) 高于教科书典型但来源如实声明。
  \item \textbf{D15}　\texttt{controls\_dc\_v\_rigid}/\texttt{dc\_p\_is\_free} 死字段（全库从未置真，Mode 4/5 刚性 Vdc 无法表达）；\texttt{k\_vdc != 0.0} 精确比较与协调侧 \texttt{|k\_vdc|>1e-9} 两容差（NaN 时结论相反）；\texttt{to\_acdc\_control\_mode} 下垂归类不对称。
  \item \textbf{D16}　\texttt{materialize\_dc\_storage} 不检查与既有 \texttt{dc.storage} 的 \texttt{.index} 冲突（\textbf{待核实}：查 validation 是否有跨容器 ID 唯一性规则）。
  \item \textbf{D17}　开关绑定拓扑推断不覆盖三绕组变压器（BFS 邻接与绑定目标均跳过）；对停运开关也生成建议；\texttt{protection\_zone\_id} 用 1000000 魔数划分命名空间未注释。
  \item \textbf{D18}　\texttt{ExternalGrid::x\_r}（$X/R$，缺省 10）与 \texttt{rx\_max/rx\_min}（$R/X$，缺省 0.1）同结构体倒数约定并存（手册已标注，命名陷阱仍在）；\texttt{PVSystem.num\_series/num\_parallel=0} 时静默零出力；VSC 损耗三口径（\texttt{eta/loss\_percent/loss\_mw}）在 CurrentBased 下同时生效易双重计损（手册已声明公式）。
  \item \textbf{D19}　\file{ac\_adjacency\_list}/\file{generators\_per\_bus} 全仓库无消费方（dead public API）；\texttt{Transformer3W} 全零容量/短路电压输入静默投影为 0 条支路无诊断。
  \item \textbf{D20}　手册覆盖缺口：\texttt{MobileStorage.grid\_forming}、Switch 上游保护开关索引字段（\texttt{upstream\_protective\_switch\_index}）及新增 \texttt{use\_signed\_series\_impedance} 三字段未进手册；\texttt{RegulatorControl}、\texttt{ThreePhaseRegulatorControl} 字段完全重复的双结构体（维护债务）。
\end{itemize}

\newpage
\section{已核查无误的数学内核（正面清单）}
\label{sec:verified}

以下核心数学内容经\textbf{解析推导 $+$ 数值验证（含链接 \file{libhacdcpf.a} 的端到端复现）}为正确（逐项行号备查）：

\begin{enumerate}
  \item \textbf{参数换算三件套}：\file{rx\_from\_vk\_vkr}（$z=\frac{v_k}{100}\frac{S_b}{S_n}$、$x=\sqrt{z^2-r^2}$，\file{network\_utils.cpp:70--82}）、\file{rx\_from\_z\_percent\_x\_over\_r}（\file{:84--97}）、\file{tap\_from\_step}（\file{:99--104}）——数值验证正确。
  \item \textbf{Transformer2W 投影}（\file{:155--234}）：电压基准折算 $(V_n/V_{base})^2$、LV 分接头 $t\to 1/t$ 且 $Z\to Zt^2$ 归算、额定变比失配 $(V_{n,H}/V_{n,L})/(V_{base,H}/V_{base,L})$、零序路径与并联台数聚合——经传递矩阵恒等式与专用回归验证正确。
  \item \textbf{Transformer3W 投影}（\file{:237--365}）：绕组对容量取 $\min$、星形合成 $Z_H=\frac{1}{2}(Z_{HM}+Z_{HL}-Z_{ML})$、星$\to$三角 $y_{ss}/(y_w y_t)$（与经典变换及网络 Kron 约化数值差 $3.6\times10^{-15}$）、三对支路 $\tau$ 表、绕组电压基准与移相角分配 $\theta_{ML}=\theta_L-\theta_M$、负阻抗腿处理——正确。
  \item \textbf{开关/断路器等值}（\file{:367--444}）：$Z_{base}=kV^2/S_b$、$Z=\max(R_c,z)$、$x=\sqrt{Z^2-R_c^2}$、理想闭合补 $10^{-6}$、$S=\sqrt{3}\,kV\cdot kA$ 热稳定——与手册 \S 4.1/\S 4.2 逐式一致。
  \item \textbf{零阻抗合并主干}（\file{:1551--1829}）：并查集、代表选择（SLACK$>$PV$>$PQ）、母线 pd/qd/gs/bs 求和迁移、组件引用整体重映射、合并前后功率守恒、白名单防误并真实短线路、\texttt{extensive\_participation} 权重和恒为 1——26+ 项断言数值复核正确。
  \item \textbf{死岛剥离与编号规范化}（\file{:1222--1548}、\file{:1839--1891}）：BFS 仅以在役支路为边、删除后重编号 1..N、BusMergeMap \textbf{组合而非覆盖}、DC 编号重复抛异常、AC/DC 域严格分离——正确。
  \item \textbf{反投影}（\file{unproject\_bus\_vector}，\file{:1896--1931}）：Intensive 广播、Extensive 按权重拆分且总量守恒、死岛置 0、缺失权重抛异常——正确。
  \item \textbf{单位换算}（\file{unit\_conversion.cpp:21--92}）：$Z_{base}=V^2/S$、串联 $r/x$ 除以 \texttt{n\_parallel}、电纳乘 \texttt{n\_parallel}、$B_{pu}=B_{S}\cdot Z_{base}$、DC 侧 $R_{base}=V^2/P$、幂等——与手册公式逐符号一致（\texttt{c\_nf} 缺陷另见 C13）。
  \item \textbf{IEC 60228/GB/T 3956 电阻查表}（\file{standard\_parameter\_library.cpp:152--173}）：20 档铜 $+$ 17 档铝与公布最大直流电阻\textbf{逐位一致}；温度修正 $\alpha_{Cu}=0.00393/\alpha_{Al}=0.00403$、电阻率回退（声明为工程假设）、载流估算（声明为筛选常数）——正确。
  \item \textbf{可靠性物化}（\file{:987--1087}）：$\lambda=\lambda'\times l$ 与下游消费闭环、MTBF/MTTR 换算（$8760/175200=0.05$ 精确）——正确（缺省 1 km 声明缺口另见 C11）。
  \item \textbf{参数库骨架}：52 条规则默认值全部在 $[min,max]$ 内且 $min\le max$；\texttt{check\_opf\_decision\_bounds} 归一化与四类诊断；效率钳制 $\eta\in(0,1]$；SOC 窗口修复（AC 侧）——正确。
  \item \textbf{ER 展开主干}：增益整定 $k_{tight}=\max(100P_N,25)$、$k_{droop}=\max(P_N/0.1,1)$ 与手册 \S 6.3 及求解器同律；$\eta_{dcdc}=1-\texttt{loss}/100$ 与 0.98 回退；AC/DC 端口按域展开、控制模式映射、先展开后 canonicalize 的编号重映射及显式 provenance record——正确。
  \item \textbf{归因层主干}（\file{result\_attribution.cpp}）：AC/DC 支路回算公式与量纲、3W pair$\to$绕组端 KCL 精确映射、shunt/储能符号（放电为正、bs 正为容性）、VSC/DCDC 终端与损耗、OPF 映射翻译（stor/ren/gen 索引空间）、\texttt{parse\_component\_index} 与 \texttt{record\_mapping} 格式匹配、域隔离及诚实 recovery 降级——正确。
  \item \textbf{头文件常量}：\file{defaults.hpp} 与结构体初值交叉一致（base\_mva=100、freq=50、vmin/vmax=0.9/1.1 等）；\file{materialize\_dc\_storage} 全部 35 字段对拷且幂等；\file{sanitize\_scaling} 与 \texttt{effective\_load\_p\_mw}（AC/DC）、AC 版 \texttt{controllable} 分支——正确。
  \item \textbf{枚举层}：Hydro 在内的枚举 to/from 成对完整（\file{BusType}、\file{ConverterMode}、\file{ACDCControlMode}、\file{LCC}、\file{DCDC}、\file{ER}、\file{FuelType}、\file{LoadModel} 等，专用回归与编译通过）。
  \item \textbf{典型参数数学}（\file{typical\_parameters.hpp}）：外网 $x=S_b/S_{sc}''$、$r=x\cdot(R/X)$、储能 $e=E_N\cdot soc$、$p_{min}=-p_{max}$ 对称——正确。
\end{enumerate}

\newpage
\section{系统性根因分析：为何既有测试未能暴露}
\label{sec:rootcause}

\begin{enumerate}
  \item \textbf{语义单一来源缺失}：字段单位/命名约定在手册、结构体、消费方三方漂移——电机 \texttt{r\_pu}（B2）、\texttt{failure\_rate}（B15）、\texttt{eol\_percent}（B20）、\texttt{mag0\_percent}（B21）、\texttt{DCBus.pd\_mw}（B22）、\texttt{self\_discharge\_pct}（B25）、\texttt{rate\_a}/\texttt{s\_max}（B24）同属一类。项目虽有 \file{defaults.hpp}/\file{effective\_capacity.hpp} 单一来源机制，但 B18 证明该机制自身契约已与实现脱节。
  \item \textbf{投影/归因``双向镜像''测试缺失}：A1、A2、B9、B10、B11 均为投影端与归因/消费端不对称（一边乘 scaling 一边不乘、一边清除一边保留、一边枚举完备一边遗漏）。任何一项 round-trip 守恒测试（投影$\to$求解$\to$归因的功率平衡闭环）都能拦截整类缺陷。
  \item \textbf{填充路径多源}：标准参数库 vs TypicalParameters vs 结构体初值三套默认值并存（B16、D14），手册 \S 1.3 的``优先级''声明未覆盖实质分歧。
  \item \textbf{手册由代码生成但已漂移}：\S 6.3 行号过期、48/52 条数、\texttt{DCBus.pd\_mw} 按旧行为撰写、ER DC 端口错误语义被照抄（A3）——手册与代码缺同步机制，审计中手册只能作为``声明意图''而非``实现事实''使用。
  \item \textbf{测试锁定退化行为}：coverage 恒真被 CHECK 固化（C17）、热态 \texttt{r\_ohm\_per\_km} 被断言固化（C10）——测试把实现抄本当预期，失去鉴别力。
  \item \textbf{条件触发型缺陷}：\texttt{n\_parallel>1}（B1/B13）、\texttt{tap\_side=1}（B4）、charger-only（A2）、DC 端口 ER（A3）、FlexibleLoad+OPF（A1）、含母线电导/相矩阵（B14）——既有测试数据全部落在未触发侧，与归档潮流审计 \file{docs/archive/audits/power\_flow\_math\_audit.tex} 发现的``算例结构盲区''根因同源。
\end{enumerate}

\section{处置建议（按优先级）}
\label{sec:recommend}

\begin{enumerate}
  \item \textbf{发布前门禁}：在与记录 pin 一致且干净的 MIPSolvers checkout 上执行正式 Release preset、完整 ctest 与可选 Ipopt/ETAP 路径。本轮 Debug 验证不能替代发布二进制认证。
  \item \textbf{回归保持}：把 \file{test\_component\_models\_math\_audit} 保留在默认测试集；后续修改 projection/graph/attribution 时，优先扩充 AC/DC 同号 ID、round-trip 守恒、dead-island 与多设备不可辨识场景，不得删除 fail-closed 断言来迁就错误输入。
  \item \textbf{文档单一来源}：将字段单位、时间基准、额定优先级和 pu 基准抽成可生成的语义登记表，消除手册行号、默认值与代码注释再次漂移的可能；B15/B20/B22/B25 的兼容口径尤其应进入 schema/API 文档。
  \item \textbf{剩余观察项}：按 D1--D20 分批处理死 API、命名重复、动态字段覆盖和导入器文档漂移；每项先确认设计意图，避免把观察项机械改成破坏兼容性的行为变更。
\end{enumerate}

\section{审计工件与可复现性}

\begin{itemize}
  \item \textbf{仓库内专用回归}：\file{tests/test\_component\_models\_math\_audit.cpp} 固化 A1--A3、投影/合并、归因、参数、单位和 validation 的关键反例；当前 14 个用例、102 个断言。
  \item \textbf{既有回归}：\file{test\_result\_attribution}、\file{test\_device\_flows}、\file{test\_typical\_parameters}、\file{test\_validation}、\file{test\_multiport\_vpp\_er}、\file{test\_three\_phase\_nr\_transformer\_source}、\file{test\_short\_circuit\_crossval}、\file{test\_short\_circuit\_iec60909\_4} 与 \file{test\_dc\_short\_circuit} 用于确认兼容面；旧错误预期仅在 C10、C17、B21 对应位置按新契约更新，IEC fixture 的电机阻抗则按 B2 从欧姆值迁移为铭牌基准 pu。
  \item \textbf{解析复算}：IEC 60228 表、温度/电阻率、2W/3W 星---三角、电压基准、并联比例、电容电纳与 ZIP 权重均保留显式公式断言或审计复算；近似/等分恢复的 recovery 等级单独核对。
  \item \textbf{构建命令}：配置与限制见第~\ref{sec:remediation}~节末尾；LaTeX 文档使用 XeLaTeX 流程编译并检查日志。Release preset 仍须在干净、匹配 pin 的兄弟依赖上补做。
  \item 本报告全部 A/B/C 处置已重新核对消费点；标注``待核实''的 D 级观察项仍需项目方确认设计意图后终裁。
\end{itemize}

\vspace{1.5em}
\noindent\textbf{报告完。}

\end{document}
